% !PRINTED 60
\weberpage{60}

\begin{center}
\textbf{§ 17. The Functions $T$.}
\end{center}

If doubly periodic functions $\varphi(u)$ of the $m$th order exist, they must be fractional functions. If we therefore set them, according to \S\,15, 10., in the form

\weq{1}{\varphi(u) = \frac{G_1(u)}{G(u)},}

where $G_1(u)$ and $G(u)$ are integral functions without a common zero point, then

\[
\frac{G_1(u+\omega_1)}{G(u+\omega_1)} = \frac{G_1(u+\omega_2)}{G(u+\omega_2)} = \frac{G_1(u)}{G(u)},
\]

and in this the three functions $G(u)$, $G(u+\omega_1)$, $G(u+\omega_2)$ have the same $m$ zero points, namely the points of discontinuity of $\varphi(u)$. Their quotients are therefore unit functions, and hence by \S\,15, 9., if $g_1(u)$, $g_2(u)$ are \gesperrt{integral functions}:

\weq{2}{\begin{aligned}
G(u+\omega_1) &= e^{g_1(u)}\,G(u),\\
G(u+\omega_2) &= e^{g_2(u)}\,G(u),
\end{aligned}}

and the function $G_1(u)$ satisfies the same equations.

If

\weq{3}{\varphi(u) = \frac{T_1(u)}{T(u)}}

is a second representation of $\varphi(u)$ of the form (1), then

\weq{4}{T(u) = e^{g(u)}\,G(u),}

where $g(u)$ is again an integral function, and if therefore, by (2),

\weq{5}{\begin{aligned}
T(u+\omega_1) &= e^{l_1(u)}\,T(u),\\
T(u+\omega_2) &= e^{l_2(u)}\,T(u)
\end{aligned}}

holds, then

\[
\begin{aligned}
g_1(u) &= l_1(u) + g(u) - g(u+\omega_1),\\
g_2(u) &= l_2(u) + g(u) - g(u+\omega_2).
\end{aligned}
\]

Everything therefore depends on our being able, under some suitable assumption about the integral functions $l_1(u)$, $l_2(u)$, to form integral functions $T(u)$ which satisfy the conditions (5) and vanish at $m$ arbitrarily given points of the period parallelogram. By the quotients of two such functions with the same $l_1(u)$, $l_2(u)$ we can then represent all doubly periodic functions $\varphi(u)$, and the more general functions $G(u)$ which accomplish the same are obtained, according to (4), by the addition of an \gesperrt{arbitrary unit factor}.
